% !TeX program = pdflatex
\documentclass[11pt]{article}
\usepackage{amsmath}
\usepackage{environ}
\usepackage{CJKutf8}

% ============ 全局三态开关 ============
\newif\ifglobalnote   % 全局值：true = 开，false = 关
\newif\ifglobalset    % 是否显式设置：true = 已设置（压制段落），false = 未设置
\globalsetfalse       % 全局未设置，段落自治

% 三个命令均用 \global\let，跨组永久生效
\newcommand{\globalnoteon}{%
	\global\let\ifglobalset\iftrue
	\global\let\ifglobalnote\iftrue}
\newcommand{\globalnoteoff}{%
	\global\let\ifglobalset\iftrue
	\global\let\ifglobalnote\iffalse}
\newcommand{\globalnoteauto}{%
	\global\let\ifglobalset\iffalse}

% ============ 段落开关常量 ============
\def\mynoteon{on}
\def\mynoteoff{off}

% ============ 中间量 ============
\newif\ifshownote

% ============ 注释段落环境 ============
% 优先级：全局已设置 > 段落选择；全局未设置时段落不写参数默认 off（不显示）
\NewEnviron{mynote}[1][off]{%
	\begingroup
	\def\mynotearg{#1}%
	\ifglobalset
	\ifglobalnote \shownotetrue \else \shownotefalse \fi
	\else
	\ifx\mynotearg\mynoteon \shownotetrue \else \shownotefalse \fi
	\fi
	\ifshownote
	\par\noindent\rule{\linewidth}{0.4pt}\par
	\BODY
	\par\noindent\rule{\linewidth}{0.4pt}\par
	\fi
	\endgroup
}

% ===== 全局状态：三选一 =====
% \globalnoteon     % ← 全局显示（所有 mynote 强制显示，无视段落参数）
\globalnoteoff  % ← 全局关闭（所有 mynote 强制隐藏，无视段落参数）
% \globalnoteauto % ← 全局未设置（段落自己决定，这是默认）
\begin{document}
	\begin{CJK*}{UTF8}{gbsn}
		
		正文。Here is an equation:
		\begin{equation}
			E = mc^2.
		\end{equation}
		
		% 全局未设置（导言区默认），段落自治
		A：[on] → 显示。
		\begin{mynote}[on]
			内容 A，含 \( \alpha+\beta \)。
		\end{mynote}
		
		B：[off] → 隐藏。
		\begin{mynote}[off]
			内容 B（不应出现）。
		\end{mynote}
		
		C：不写 → 默认 off，隐藏。
		\begin{mynote}
			内容 C（不应出现）。
		\end{mynote}
		
		% 中途切换全局：无论写在 CJK* 内外、组内组外都永久生效
		\globalnoteoff
		D：[on] 但全局关 → 隐藏。
		\begin{mynote}[on]
			内容 D（不应出现）。
		\end{mynote}
		
		\globalnoteon
		E：[off] 但全局开 → 显示。多行公式：
		\begin{mynote}[off]
			内容 E。
			\begin{align}
				a &= b + c, \\
				d &= e + f.
			\end{align}
		\end{mynote}
		
		\globalnoteauto   % 恢复段落自治；全文只需这一处（或干脆不写）
		F：[off] → 隐藏。
		\begin{mynote}[off]
			内容 F（不应出现）。
		\end{mynote}
		
		G：[on] → 显示。
		\begin{mynote}[on]
			内容 G，含中文与公式 \( \gamma \)。
		\end{mynote}
		
	\end{CJK*}
\end{document}